\documentclass[book.tex]{subfiles}

\begin{document}
\chapter{Variational Quantum Algorithms}
\label{chap:variation_quantum_algorithms}
\noindent\rule{11cm}{0.4pt} \\
\textbf{Prerequisites:} \autoref{chap:quantum_circuits} \\
\textbf{Difficulty Level:} ** \\
\noindent\rule{11cm}{0.4pt}
\vspace{5mm}

A variational quantum algorithm is a quantum circuit parameterized by a given set of parameters $\theta$. Variational quantum algorithms can be used to combine classical learning techniques with quantum circuits. The main motivation for using VQA algorithms is that noisy intermediate scale qubit (NISQ) devices have high noise, high error, limited connectivity and cannot yet make use of quantum error correction. Intuitively, the parameters act as a mechanism for controlling the error on these devices.

VQA algorithms off-board parameterization to classical machine learning algorithms and use the minimal viable quantum circuit for the task at hand.

\section{Basic Formulation}

A variational quantum problem is set up much like a standard machine learning problem. A loss, denoted $C$ is optimized to find an optimal set of parameters. Optimization is performed with standard machine learning algorithms.
\begin{align}
    \theta^* = \arg\min_{\theta} C(\theta)
\end{align}

The variation is that the loss function $C$ is implemented by a quantum circuit. The general quantum loss function is expressed in the form

\begin{align}
    C(\theta) &= \sum_k f_k(\mathrm{Tr}[O_kU(\theta)\rho_kU^{\dagger}(\theta)])
\end{align}

Here $U(\theta)$ is a parameterized unitary transformation, $\rho_k$ are the input training samples, $O_k$ are the observables, and $f_k$ are arbitrary functions. 

The form of a parameterized unitary $U(\theta)$ is called its ansatz. The generic $U(\theta)$ can be broken down as a sequence of simpler unitary transformations

\begin{align}
    U(\theta) = U_L(\theta_L) \dotsc U_1(\theta_1)
\end{align}

and each $U_i(\theta)$ can be further decomposed as

\begin{align}
    U_i(\theta) &= \prod_j e^{-i\theta_jH_j} W_j
\end{align}
where the $H_j$ are Hermitian operators and the $W_j$ are unparameterized unitary transformations. At a high level, this decomposition is similar to decompositions seen in deep networks, which sequences of chained transformations. An ansatz then corresponds roughly to a family of architectures (like convolutional networks or transformers).

\section{The Parameter Shift Rule}

For classical differentiable programs, we use automatic differentiation methods to find gradients. But $C$ is a quantum loss; how can we find the gradient of $C$ with respect to its parameters? The parameter shift rule is a variant of the finite difference method which suffices to take derivatives.

Suppose that $f_k$ is the identity function, and that $\theta_l$ parameterizes unitary $e^{i\theta_l \sigma_l}$ where $\sigma$ is a Pauli operator. Then the parameter shift rule is given by

\begin{align}
    \frac{\partial C}{\partial \theta_l} &= \sum_k \frac{1}{\sin 2 \alpha } \left ( C(\theta_l + \frac{\pi}{4l}) - C(\theta_l - \frac{\pi}{4l}) \right )
\end{align}

%\section{Variational Quantum Eigensolver}

%\textcolor{red}{TODO: Add description of variational quantum eigensolver}


%\section{Quantum Neural Networks}

%\textcolor{red}{TODO: Universal quantum circuit simulation?}


\end{document} 
